% theme: changes_of_state
\MajorQuestion{物質の状態変化}
\SetRowSpaceQanda

\Instruction{次の問いに答えなさい。}
\QQ{温度によって物質の状態が変化することを\NamedBlank{state_change}という。}{状態変化}
\QQ{物質をつくっている非常に小さなものを\NamedBlank{particle}という。}{粒子}
\QQ{液体を沸騰させ、出てきた気体を冷やして再び液体として取り出す方法を\NamedBlank{distillation}という。}{蒸留}
\QQ{固体・液体・気体の3つの状態をまとめて何というか。}{物質の三態}
\QQ{\RefBlank{particle}が規則正しく並び、その場で振動している状態を\NamedBlank{solid}という。}{固体}

\Instruction{\strut}
\QQ{\RefBlank{particle}が近くに集まり、互いに動き回っている状態を\NamedBlank{liquid}という。}{液体}
\QQ{\RefBlank{particle}が離れて存在し、自由に飛び回っている状態を\NamedBlank{gas}という。}{気体}
\QQ{固体から直接気体に変化する\RefBlank{state_change}を何というか。}{昇華}
\QQ{\RefBlank{distillation}では、液体どうしの何の違いを利用して物質を分けるか。}{沸点の違い}
\QQ{液体を加熱するとき、急な沸騰を防ぐために入れるものを何というか。}{沸騰石}

\Instruction{\strut}
\QQ{\RefBlank{solid}がとけて\RefBlank{liquid}に変化するときの温度を\NamedBlank{melting_point}という。}{融点}
\QQ{一般に、固体・液体・気体のうち体積が最も大きい状態はどれか。}{気体}
\QQ{水が\RefBlank{solid}の氷になるとき、体積は大きくなるか小さくなるか。}{大きくなる}
\QQ{水とエタノールの混合物を\RefBlank{distillation}したとき、はじめに得られる液体により多く含まれる物質はどちらか。}{エタノール}
\QQ{\RefBlank{liquid}が沸騰して\RefBlank{gas}に変化するときの温度を\NamedBlank{boiling_point}という。}{沸点}

\Instruction{\strut}
\QQ{物質が\RefBlank{state_change}するとき、変化しない量は質量と体積のどちらか。}{質量}
\QQ{\RefBlank{liquid}が液面から\RefBlank{gas}に変化することを何というか。}{蒸発}
\QQ{\RefBlank{state_change}中も温度が変化するのは純物質と混合物のどちらか。}{混合物}
\QQ{\RefBlank{liquid}の内部からも\RefBlank{gas}に変化することを何というか。}{沸騰}
\QQ{一定の\RefBlank{melting_point}や\RefBlank{boiling_point}を示すのは、純物質と混合物のどちらか。}{純物質}
